\documentclass{exam}

\newif\ifA
\newif\ifkey
\newif\ifprintselected   % required by exam_config's \ansbox

%%%%%%%%%%%%%%%%%%%%%%%%%%%%%%%%%%%%%%%%%%%%%%%%%%%%%%
%%% SET THESE IF VARIABLES
%%%%%%%%%%%%%%%%%%%%%%%%%%%%%%%%%%%%%%%%%%%%%%%%%%%%%%
% Lab Num -- Module 3's six lab periods follow L04 (Lab 1) and the two
% Module 2 measurement labs, so L25 is Lab 8. Override with \def\labNumIn{}
% if the course sequence renumbers.
\ifdefined\labNumIn
	\def\labNum{\labNumIn}
\else
	\def\labNum{8}
\fi

% is this the key or not?
\ifdefined\iskey
	\ifnum\iskey=1
		\keytrue
	\else
		\keyfalse
	\fi
\else
	%% Pick one manually
	\keyfalse % if this is NOT the key
%	\keytrue     % if this IS the key
\fi

% set up the header and footer - change for every term
\ifdefined\thisTermIn
	\def\thisterm{\thisTermIn}
\else
	\def\thisterm{Fall 2026}
\fi
%
\def\thisclass{ECE 444}

%%%%%%%%%%%%%%%%%%%%%%%%%%%%%%%%%%%%%%%%%%%%%%%%%%%%%%
%%% END SET THESE IF VARIABLES
%%%%%%%%%%%%%%%%%%%%%%%%%%%%%%%%%%%%%%%%%%%%%%%%%%%%%%
% packages and shortcuts
\input{myPackages_gr}
\usepackage[hidelinks, linkcolor=red]{hyperref}
\input{myShortcuts}
\setlength{\parindent}{0pt}
\setlength{\parskip}{6pt}

%% exam class customization
\input{exam_config}
\printselectedfalse      % full worked solutions in the key, not answers-only
\CorrectChoiceEmphasis{\color{red}}
\SolutionEmphasis{\color{red}}
\renewcommand\questionlabel{\textbf{\thequestion.}}
\renewcommand{\thepartno}{\alph{partno}}
\renewcommand{\partlabel}{(\thepartno)}

% Fillable table cell: blank in the student copy, the expected value in red on
% the KEY. Fixed-width centered columns keep the blank cells writable.
\newcolumntype{K}[1]{>{\centering\arraybackslash}p{#1}}
\newcommand{\kv}[1]{\ifkey{\color{red}#1}\fi\rule[-0.10in]{0pt}{0.34in}}

% print the answers if this is the key
\ifkey
	\printanswers % if this IS the key
\else
	\noprintanswers  %if this is NOT the key
\fi

% print the header and footer
% need to print key on top if this is the key
\ifkey
	\firstpageheader{}{\color{red}{ KEY}}{}
	\runningheader{\thisclass\hspace{1pt} -- \thisterm }{Lab \#\labNum\hspace{1pt}\color{red}{ KEY}}{\thepage\hspace{1pt} of\hspace{1pt} \numpages}
	\firstpagefooter{}{}{}
	\runningfooter{}{}{}
\else
	\firstpageheader{}{}{}
	\runningheader{\thisclass \hspace{1pt} -- \thisterm }{Lab \#\labNum\hspace{1pt}}{\thepage\hspace{1pt} of\hspace{1pt} \numpages}
	\firstpagefooter{}{}{}
	\runningfooter{}{}{}
\fi

\runningheadrule

\begin{document}
\pagestyle{headandfoot}   % exam class


%%%%%%%%%%%%%%%%%%%%%%%%%%%%%%%%%%%%%%%%%%%%%%%%%%%%%%
% title block
%%%%%%%%%%%%%%%%%%%%%%%%%%%%%%%%%%%%%%%%%%%%%%%%%%%%%%

\begin{center}
USAF ACADEMY \\
DEPARTMENT OF ELECTRICAL AND COMPUTER ENGINEERING \\
LAB \#\labNum\ (Lesson 25) - Tapering Lab \\
\textbf{Lab Sheet} \\
\thisterm \\[4mm]
\end{center}

\vspace{-4pt}
\textbf{Name:} \rule[-2pt]{2.6in}{0.4pt} \hfill
\textbf{Section:} \rule[-2pt]{0.9in}{0.4pt} \hfill
\textbf{Date:} \rule[-2pt]{1.1in}{0.4pt}

\vspace{4pt}

\fullwidth{\textbf{Learning Objective 3.7: } Apply amplitude tapering to control
sidelobe level, and predict the resulting pattern trade-off.}

\vspace{2pt}

This sheet is the turn-in document for the Lesson 25 lab. Predict each row
before you sweep it, and keep the peak drop and the taper efficiency in separate
columns throughout. The lab is a team assignment, so one sheet per team.

\begin{questions}

%%%%%%%%%%%%%%%%%%%%%%%%%%%%%%%%%%%%%%%%%%%%%%%%%%%
\question \textbf{The taper table, part 1: the element gains.} Read the eight
slider values off Element Gains for each preset and record them as percentages.

\begin{center}
\small
\setlength{\tabcolsep}{4pt}
\renewcommand{\arraystretch}{1.25}
\begin{tabular}{| l | K{0.5in} | K{0.5in} | K{0.5in} | K{0.5in} | K{0.5in} | K{0.5in} | K{0.5in} | K{0.5in} |}
\hline
\textbf{Preset} & \textbf{Rx1} & \textbf{Rx2} & \textbf{Rx3} & \textbf{Rx4} & \textbf{Rx5} & \textbf{Rx6} & \textbf{Rx7} & \textbf{Rx8} \\\hline\hline
Uniform & \kv{100} & \kv{100} & \kv{100} & \kv{100} & \kv{100} & \kv{100} & \kv{100} & \kv{100} \\\hline
Hann & \kv{12} & \kv{43} & \kv{77} & \kv{100} & \kv{100} & \kv{77} & \kv{43} & \kv{12} \\\hline
Blackman & \kv{6} & \kv{27} & \kv{66} & \kv{100} & \kv{100} & \kv{66} & \kv{27} & \kv{6} \\\hline
Chebyshev & \kv{4} & \kv{23} & \kv{62} & \kv{100} & \kv{100} & \kv{62} & \kv{23} & \kv{4} \\\hline
\end{tabular}
\end{center}

\textbf{The taper table, part 2: what the sweep reads.} Predict each row before
you press Start. The peak drop is measured against the frozen uniform trace, and
$\eta_t = (\sum a_n)^2 / (N \sum a_n^2)$ is computed rather than measured.

\begin{center}
\footnotesize
\setlength{\tabcolsep}{3pt}
\renewcommand{\arraystretch}{1.25}
\begin{tabular}{| l | K{0.6in} | K{0.65in} | K{0.6in} | K{0.95in} | K{0.9in} | K{0.45in} | K{0.55in} |}
\hline
\textbf{Preset} & \textbf{HPBW pred} & \textbf{HPBW meas} & \textbf{Drop pred} & \textbf{Drop meas} & \textbf{First sidelobe} & $\boldsymbol{\eta_t}$ & $\boldsymbol{\eta_t}$ \textbf{(dB)} \\\hline\hline
Uniform & \kv{$13.1\mydeg$} & \kv{$\approx 13\mydeg$} & \kv{$0$ dB} & \kv{$0$ dB} & \kv{$-11$ to $-13$ dBc} & \kv{1.00} & \kv{$0$} \\\hline
Hann & \kv{$19.5\mydeg$} & \kv{$19$ to $20\mydeg$} & \kv{$-4.7$ dB} & \kv{$-4.7 \pm 0.5$ dB} & \kv{below the noise floor} & \kv{0.75} & \kv{$-1.2$} \\\hline
Blackman & \kv{$23.1\mydeg$} & \kv{$22$ to $24\mydeg$} & \kv{$-6.1$ dB} & \kv{$-6.1 \pm 0.5$ dB} & \kv{below the noise floor} & \kv{0.66} & \kv{$-1.8$} \\\hline
Chebyshev & \kv{$24.3\mydeg$} & \kv{$23$ to $25\mydeg$} & \kv{$-6.5$ dB} & \kv{$-6.5 \pm 0.5$ dB} & \kv{below the noise floor} & \kv{0.62} & \kv{$-2.1$} \\\hline
\end{tabular}
\end{center}

\begin{solution}
Work the Hann row through as the model. With
$a_n = 0.12,\ 0.43,\ 0.77,\ 1.00,\ 1.00,\ 0.77,\ 0.43,\ 0.12$,
\eq{
	\sum a_n &= 4.64, \qquad \sum a_n^2 = 3.584 \\
	\text{peak drop} &= 20\mylog{\frac{4.64}{8}} = 20\mylog{0.580} = -4.7 \text{ dB} \\
	\eta_t &= \frac{(4.64)^2}{8(3.584)} = \frac{21.53}{28.67} = 0.751 \quad \rightarrow \quad -1.2 \text{ dB}
}
The same two sums give $-6.1$ dB and $0.655$ for Blackman ($\sum a_n = 3.98$,
$\sum a_n^2 = 3.024$), and $-6.5$ dB and $0.621$ for Chebyshev
($\sum a_n = 3.78$, $\sum a_n^2 = 2.878$). Measured beamwidths land within a
degree or so of the prediction, because the sweep steps in $2.8125\mydeg$
increments and the 3 dB crossings fall between samples. Every tapered sidelobe
sits below the sweep's floor, which is about 23 dB under the uniform peak, so
the entry in that column is the phrase and not a number read off the grass.
\end{solution}

%%%%%%%%%%%%%%%%%%%%%%%%%%%%%%%%%%%%%%%%%%%%%%%%%%%
\newpage
\question \textbf{Your custom taper.} Record the eight gain values you settled
on and the measured result. The specification was a half-power beamwidth no
wider than $17\mydeg$ with the first sidelobe below $-20$ dBc.

\begin{center}
\small
\setlength{\tabcolsep}{4pt}
\renewcommand{\arraystretch}{1.25}
\begin{tabular}{| l | K{0.5in} | K{0.5in} | K{0.5in} | K{0.5in} | K{0.5in} | K{0.5in} | K{0.5in} | K{0.5in} |}
\hline
\textbf{Taper} & \textbf{Rx1} & \textbf{Rx2} & \textbf{Rx3} & \textbf{Rx4} & \textbf{Rx5} & \textbf{Rx6} & \textbf{Rx7} & \textbf{Rx8} \\\hline\hline
Custom & \kv{$\approx 45$} & \kv{$\approx 65$} & \kv{$\approx 85$} & \kv{100} & \kv{100} & \kv{$\approx 85$} & \kv{$\approx 65$} & \kv{$\approx 45$} \\\hline
\end{tabular}
\end{center}

\vspace{4pt}
\begin{center}
Measured HPBW = \ansbox[1.4in]{$\le 17\mydeg$} %
Measured first sidelobe = \ansbox[1.8in]{below $-20$ dBc}
\end{center}
\vspace{4pt}

In one sentence, how did you arrive at these values?

\begin{solutionorlines}[1.2in]
Any symmetric taper that meets the specification is acceptable, so there is no
single set of eight numbers to look for. Grade the arithmetic instead: the peak
drop reported has to equal $20\mylog{\sum a_n / 8}$ for the gains written down.
A taper with the end elements near 45 percent, stepping smoothly up to 100
percent in the middle, reads roughly $15\mydeg$ of beamwidth and a peak drop
near $-3.3$ dB, with the first sidelobe just reaching the floor. The last line
should describe two or three iterations: set the sliders, sweep, read the HPBW
and the sidelobe level, and adjust.
\end{solutionorlines}

%%%%%%%%%%%%%%%%%%%%%%%%%%%%%%%%%%%%%%%%%%%%%%%%%%%
\newpage
\question \textbf{Written answers.} A short paragraph each.

\begin{parts}

	\part The tapered sidelobes disappear from the plot, but the beamwidth change
	is easy to see and easy to measure. Explain why the measurement gives a good
	number for one and no number at all for the other.

	\begin{solutionorlines}[2.2in]
	The beamwidth is read near the top of the trace, 3 dB below a peak that stands
	about 23 dB above the noise floor, so both crossings are well inside the
	dynamic range of the measurement. The tapered sidelobes are pushed more than
	17 dB below that peak, which puts them under the floor, and anything below the
	floor reads as noise rather than as pattern. The instrument is not measuring
	the sidelobe at all in those rows, so a number taken from the trace would
	report the receiver's noise. The correct entry is that the sidelobe is below
	the noise floor.
	\end{solutionorlines}

	\part Explain why the peak drop you measured is not the directivity the array
	lost, and what each number would be used for.

	\begin{solutionorlines}[2.6in]
	The peak drop is the coherent receive-voltage loss,
	$20\mylog{\sum a_n / N}$, read off a trace that is normalized to the uniform
	peak. The tapered array does collect less signal voltage from the source, and
	that is the $4.7$ dB the Hann trace falls. It also responds to less of
	everything arriving from off-boresight, and directivity is the ratio of
	on-axis intensity to the average over all angles, so that ratio falls only by
	the taper efficiency, $\eta_t = 0.75$, or $1.2$ dB. Re-normalizing each trace
	to its own peak would remove the $4.7$ dB from the plot entirely and leave the
	beam shape unchanged. Use the peak drop to predict what the plot will show and
	to check that the intended gains were loaded, and use $\eta_t$ in a link or
	system budget, where debiting $4.7$ dB would overstate the loss by a factor of
	two in power.
	\end{solutionorlines}

\end{parts}

\end{questions}

\vspace{1.0em}
\noindent\textbf{Documentation:} \rule[-2pt]{0.6\linewidth}{0.4pt}

\end{document}
